An important class of enhanced sampling methods is to modify the topology of the underlying PES to ease the barrier crossing, as represented by metadynamics pioneered by \citet{laio2002escaping} [\cref{fig:sampling_methods}(i)].  In order to cross the energy barrier that separates the two energy basins in MD, a time and position bias potential $V(\bm{R},t)$ is progressively added to the underlying energy landscape to `fill up' the basins, until some point $t_{final}$ where the energy basins become completely `flat' and the system can move without any energy barrier between the two basins.  Because of the dynamic nature of the MD simulations, they not only provide information on the potential energy but also allow direct access to the system's free energy that accounts for both the vibrational and configurational degrees of freedom. This leads to the unique advantage of applying metadynamics in PES sampling whereby, upon convergence, the negative of the total bias potential added during the simulation gives the system's free energy landscape, $i.e.$ $F(\bm{R})=-V(\bm{R},t_{final})$.

Practically, the bias potential in  metadynamics  is added to the so-called \emph{collective variables} (CVs), which are  (or derived from) a subset of all the atomic coordinates in the system. They are a `focused' set of degrees of freedom which usually define the slowest parts of the system dynamics that are of interest. This implies that when the system is traversing the space spanned by the CVs, it needs to cross an energy barrier that is significantly larger than $K_B T$. 

Here, we will provide some examples of the typical choices of CVs for applying metadynamics in molecular and materials simulations.  For chemical reactions involving small molecules, using interatomic distances as CVs is the simplest choice and usually leads to meaningful chemical insights, as seen in the examples of studying the autoionisation of water \cite{liu2023mechanistic} as well as the kinetics of oxygen reduction reactions at the air/water interface for catalyst designs.\cite{yang2023neural} Bond angles consistute a set of CVs for examining the conformational dynamics of small molecules such as peptitdes,\cite{mccarty2017variational}  whereas the radius of gyration is typically selected for simulating the folding dynamics of polypeptides and proteins.\cite{vymetal2011gyration} For solid-state systems, the Steinhardt order parameters\cite{steinhardt1983bond} are often used as indicators for the emergence of cubic and hexagonal orderings from disordered structures by measuring the angular distributions of particles around a given structural centre. Both the standard form of the order parameters and their extension have been applied for simulating the nucleation of ice\cite{quigley2008metadynamics} and molecular crystals\cite{giberti2015metadynamics} with metadynamics.  A closely related type of CV is the coordination numbers\cite{iacovella2007efficient} which can be used for monitoring the changes of local coordination environments during the simulation. The lattice vectors\cite{martovnak2003predicting} are an effective set of CVs for simulating the phase transformations of crystalline materials, which comes naturally from the ideas behind the Parrinello-Rahman barostat.\cite{parrinello1980crystal} 

Identifying a set of suitable CVs is critical in metadynamics simulations. Traditionally, CVs are chosen based on chemical intuitions to identify the slowest degrees of freedom along the pathways of chemical reactions or phase transitions, which may not always be a trivial task for complex systems such as proteins and solids with more than a few atoms in the unit cells. The dimensionality of the CVs also needs to be controlled to lower the computational complexities and costs while ensuring all the relevant slow modes in the system have been sufficiently sampled. In this sense, the CVs identified for a particular system are usually not transferable to another. Finally, as the usage of CVs essentially reduces the dimensionality of the sampling problem, it leads to difficulties in judging whether hidden barriers or other important dynamics have been missed in the sampling, particularly for complex material systems. In this regard, the wide adaptation of machine-learning in atomistic simulations has opened new pathway to identify CVs that is solely based on the intrinsic dynamics of the systems itself, for example, using dimensional reduction techniques such as principle component analysis\cite{sittel2018perspective} to identify the slowest degree of freedom from existing MD trajectories. 
